\clearpage
\begingroup
\let\clearpage\relax
\let\cleardoublepage\relax
\vspace*{\fill}
\chapter{运输层}
\begin{center}
运输层为不同主机上的「进程」提供逻辑通信。本章从多路复用/分用出发，讲清 UDP、可靠数据传输、TCP 连接管理与流量/拥塞控制。
\end{center}
\vspace*{\fill}
\endgroup
\clearpage

\section{运输层概述与多路复用/分用}

\subsection{从主机到进程：逻辑通信}

网络层把分组从一台主机送到另一台主机（\emph{主机到主机}交付）。但真正收发数据的不是主机，而是主机上运行的一个个\emph{进程}。运输层在网络层服务的基础上，把交付粒度从「主机」细化到「进程」，为运行在不同主机上的应用进程之间提供\emph{逻辑通信 (logical communication)}：从应用的角度看，两个进程仿佛由一条直接的管道相连，而无需关心下面用的是哪条链路、经过了哪些路由器。

运输层协议运行在端系统的\emph{操作系统内核}中，其数据单元常称为\emph{报文段 (segment)}。因特网提供两种运输层协议：

\begin{itemize}
    \item \textbf{UDP (User Datagram Protocol)}：无连接、不可靠、尽力而为。
    \item \textbf{TCP (Transmission Control Protocol)}：面向连接、可靠、字节流、全双工。
\end{itemize}

\subsection{多路复用与多路分用}

一台主机上同时运行着浏览器、邮件客户端、下载工具等多个进程，它们都在收发数据。运输层必须解决两个方向的问题：

\begin{itemize}
    \item \textbf{多路复用 (multiplexing)}：发送端从多个 socket 收集数据，为每块数据加上运输层首部（其中含端口号等标识），再交给网络层发送。
    \item \textbf{多路分用 (demultiplexing)}：接收端依据运输层首部中的字段，把收到的报文段交付到正确的 socket。
\end{itemize}

分用的关键是\emph{端口号}。端口号是 $16$ 位整数，取值 $0\sim65535$：$0\sim1023$ 为\emph{周知端口}（如 HTTP 用 $80$、DNS 用 $53$），$1024\sim49151$ 为\emph{登记端口}，$49152\sim65535$ 为\emph{临时端口}（客户机动态申请）。

\subsection{UDP 二元组分用与 TCP 四元组分用}

两种协议的分用依据不同，这是本章第一个容易混淆的点。

\begin{table}[H]
\centering
\caption{UDP 与 TCP 的分用依据对比}
\begin{tabulary}{\textwidth}{CC}
\toprule
协议 & 分用元组 \\
\midrule
UDP & (目的 IP 地址, 目的端口号) \\
TCP & (源 IP 地址, 源端口号, 目的 IP 地址, 目的端口号) \\
\bottomrule
\end{tabulary}
\label{tab:demux}
\end{table}

\textbf{为什么 UDP 只需二元组？} UDP 是\emph{无连接、无状态}的。当应用调用 \texttt{socket()} 创建 UDP 套接字并 \texttt{bind} 到本地某端口后，该 socket 只与「本地端口」绑定，不关心对端是谁。因此，凡是目的 IP 为本机、目的端口等于该端口的所有报文段，无论来自哪个源 IP、哪个源端口，都被交付到\emph{同一个} socket。分用只需查 (目的 IP, 目的端口) 这两项。

\textbf{为什么 TCP 需要四元组？} TCP 是\emph{面向连接、有状态}的。服务器用一个\emph{监听套接字}（欢迎套接字）在端口（如 $80$）上等待连接。每当一个连接建立，内核会为该连接\emph{新建一个专属套接字}，并用四元组 (源 IP, 源端口, 目的 IP, 目的端口) 唯一标识它。这样，即使多个客户机用不同源端口、甚至相同源 IP 的不同源端口同时访问服务器的 $80$ 端口，内核也能把它们区分开来，各自维护独立的连接状态（序号、窗口、缓冲等）。若只用二元组，多个连接就会共用一个 socket、无法区分，也不可能各自维护状态。

\begin{figure}[H]
\centering
\begin{tikzpicture}[font=\small, >=latex, thick,
    host/.style={draw, rounded corners, minimum width=2.3cm, minimum height=0.9cm, align=center, fill=blue!5},
    sock/.style={draw, minimum width=2.4cm, minimum height=0.65cm, align=center, fill=orange!12},
    pkt/.style={->, thick}]
    % (a) UDP
    \begin{scope}
        \node[host] (a1) at (0,1.8) {主机 A};
        \node[host] (a2) at (0,0.2) {主机 B};
        \node[sock] (au) at (4.6,1.0) {UDP socket :53};
        \draw[pkt] (a1.east) -- node[above, font=\scriptsize]{$:3000\to:53$} (au.west);
        \draw[pkt] (a2.east) -- node[below, font=\scriptsize]{$:4000\to:53$} (au.west);
        \node[font=\scriptsize, align=center] at (2.3,-0.7) {(a) UDP：不同来源\\汇入同一 socket};
    \end{scope}
    % (b) TCP
    \begin{scope}[shift={(8.2, 0)}]
        \node[host] (b1) at (0,2.2) {主机 A};
        \node[host] (b2) at (0,0.2) {主机 B};
        \node[sock] (bs1) at (4.6,2.2) {TCP socket \#1 :80};
        \node[sock] (bs2) at (4.6,0.2) {TCP socket \#2 :80};
        \draw[pkt] (b1.east) -- (bs1.west);
        \draw[pkt] (b2.east) -- (bs2.west);
        \node[font=\scriptsize, align=center] at (2.3,-0.7) {(b) TCP：不同四元组\\映射到不同 socket};
    \end{scope}
\end{tikzpicture}
\caption{多路分用：UDP 二元组 vs TCP 四元组}
\label{fig:demux}
\end{figure}

\section{UDP：用户数据报协议}

\subsection{UDP 首部逐字段}

UDP 只做「复用/分用 + 少量差错检测」两件事，因此首部极简，固定 \textbf{8 字节}，共四个字段。

\begin{figure}[H]
\centering
\begin{tikzpicture}[font=\small]
    \matrix (h) [matrix of nodes,
        nodes={draw, minimum width=3.1cm, minimum height=0.95cm, anchor=center, align=center},
        column sep=-\pgflinewidth, row sep=-\pgflinewidth] {
        16 位源端口号 & 16 位目的端口号 \\
        16 位 UDP 长度 & 16 位检验和 \\
    };
    \node[font=\scriptsize, anchor=east] at ($(h.west)+(-0.15,0)$) {32 位};
\end{tikzpicture}
\caption{UDP 首部（8 字节）}
\label{fig:udp_header}
\end{figure}

\begin{tabulary}{\textwidth}{CC}
\toprule
字段 & 含义与要点 \\
\midrule
源端口号 & 发送进程端口，接收方回信时用作目的端口；可为 $0$ 表示不需要回复 \\
目的端口号 & 接收进程端口，是分用的核心依据 \\
UDP 长度 & 首部 $+$ 数据的总字节数，最小值 $8$（仅首部，无数据） \\
检验和 & 覆盖「伪首部 $+$ UDP 首部 $+$ UDP 数据」，用于差错检测（见下） \\
\bottomrule
\end{tabulary}

UDP 的\emph{长度}字段只含本 UDP 数据报（首部加数据），不含 IP 首部；而 IP 总长度含 IP 首部，二者不同，这是常见笔误点。

\subsection{UDP 检验和：16 位反码求和算例}

UDP 检验和的计算规则：

\begin{enumerate}
    \item 把检验和字段先置为 $0$，并在数据报前临时拼上一个 $12$ 字节的\emph{伪首部}：源 IP ($32$ 位)、目的 IP ($32$ 位)、全零 ($8$ 位)、协议号 ($8$ 位，UDP 为 $17$)、UDP 长度 ($16$ 位)。
    \item 把整个范围按 $16$ 位分组，做\emph{反码和}：逐组相加，凡产生最高位进位（第 $17$ 位）就把进位回卷加到最低位（end-around carry）。
    \item 把最终的 $16$ 位和取反码，即得检验和；若结果为全 $0$，发送时写全 $1$（避免与「不计算」混淆）。
    \item 接收方把包括检验和在内的全部 $16$ 位字再做一次反码和，结果应为 \texttt{0xFFFF}。
\end{enumerate}

\textbf{完整算例}。为便于手算，设三组 $16$ 位字为：

\[
\begin{array}{c|c}
\text{内容} & \text{16 位值} \\
\hline
w_1 & 0110\,0100\,1001\,0001 = \texttt{0x6491} \\
w_2 & 0101\,0101\,0101\,0101 = \texttt{0x5555} \\
w_3 & 1000\,1000\,1000\,1000 = \texttt{0x8888} \\
\end{array}
\]

按步计算：

\begin{align*}
w_1 + w_2 &= \texttt{0x6491} + \texttt{0x5555} = \texttt{0xB9E6} \\
\texttt{0xB9E6} + w_3 &= \texttt{0xB9E6} + \texttt{0x8888} = \texttt{0x1\,426E} \quad(\text{产生进位}) \\
\text{回卷进位} &: \texttt{0x426E} + 1 = \texttt{0x426F} \\
\text{取反码} &: \overline{\texttt{0x426F}} = \texttt{0xBD90}
\end{align*}

因此发送的检验和为 \texttt{0xBD90}。接收方验证：

\[
w_1 + w_2 + w_3 + \text{检验和}
= \texttt{0xB9E6} + \texttt{0x8888} + \texttt{0xBD90}
= \texttt{0x1\,426E} + \texttt{0xBD90}
\;\Longrightarrow\; \texttt{0x426F} + \texttt{0xBD90} = \texttt{0xFFFF}
\]

结果全 $1$，校验通过。反码和能检测所有单比特错误以及绝大多数多比特错误，但无法发现「两组字同一位同时翻转」这类使和不变的错误。

\subsection{UDP 的适用场景}

UDP 牺牲可靠性换来低时延、低开销与灵活控制：

\begin{itemize}
    \item \textbf{无握手时延}：无需连接建立，适合 DNS 这类一次问答的应用。
    \item \textbf{无连接状态}：服务器不必维护连接表，可支持海量并发客户端。
    \item \textbf{首部开销小}：仅 $8$ 字节，适合小报文。
    \item \textbf{应用层可控}：不做拥塞控制，发送速率由应用决定，适合实时音视频、VoIP、在线游戏（可容忍丢包但不能容忍重传时延）。
    \item \textbf{支持多播/广播}：TCP 不支持一对一以外的传输。
\end{itemize}

典型使用 UDP 的协议有 DNS、DHCP、TFTP、SNMP、RIP、NTP。若应用既要 UDP 的速度又要可靠性，可在\emph{应用层}自行实现确认与重传（如 QUIC 在 UDP 之上重建可靠与拥塞控制）。

\section{可靠数据传输原理}

\subsection{从理想信道到有差错信道}

可靠数据传输（reliable data transfer, rdt）的核心是：在底层信道可能\emph{丢包}或\emph{比特出错}的情况下，向上层提供「不丢、不重、不乱序」的交付。逐步引入的机制包括：

\begin{itemize}
    \item \textbf{差错检测}：检验和，使接收方能发现错误。
    \item \textbf{确认 (ACK)}：接收方告知「已正确收到」。
    \item \textbf{序号 (seq)}：区分新分组与重传/重复分组。
    \item \textbf{超时重传 (timeout + retransmit)}：发送方在合理时间内未收到 ACK 就重发。
    \item \textbf{窗口 (window)}：允许连续发送多个未确认分组，提高利用率。
\end{itemize}

\subsection{停等协议及其利用率推导}

停等 (stop-and-wait) 协议：发送方每发一个分组，就停下来等待 ACK，收到后再发下一个。

设分组长度 $L$ 比特，链路带宽 $R$ 比特/秒，往返时延为 $RTT$（忽略 ACK 传输与处理时间）。一次「发送---确认」周期由三部分组成：发送时延 $L/R$、传播往返 $RTT$、再次发送前的等待。一个周期的总时间为 $RTT + L/R$，其中\emph{真正用于发送有用数据}的时间只有 $L/R$。于是链路的\emph{利用率}为

\begin{equation}
    U_{\text{stop-and-wait}} = \frac{L/R}{RTT + L/R}.
    \label{eq:stopwait}
\end{equation}

\textbf{算例}。设 $L = 8000$ 比特，$R = 1$ Mbps，$RTT = 30$ ms。则发送时延 $L/R = 8000/10^6 = 8$ ms，

\[
U = \frac{8}{30 + 8} = \frac{8}{38} \approx 0.21.
\]

链路带宽仅被利用了约 $21\%$。$RTT$ 越大（如跨洲链路）利用率越低，这正是停等协议的致命缺陷。

\subsection{流水线与滑动窗口}

为了填满「等待 ACK」的空闲，允许发送方连续发送多个未被确认的分组，称为\emph{流水线 (pipelining)}。若在途分组数最多为 $N$（窗口大小），在\emph{无丢包、无出错}时利用率升为

\begin{equation}
    U_{\text{window}} = \frac{N \cdot L/R}{RTT + L/R},
    \qquad N \ge \frac{RTT + L/R}{L/R} \;\Rightarrow\; U = 1.
    \label{eq:window}
\end{equation}

\textbf{算例}。沿用上例，$L/R = 8$ ms、$RTT = 30$ ms，分母 $38$ ms。取 $N = 3$：$U = 24/38 \approx 0.63$；取 $N = 5$：$N L/R = 40 > 38$，利用率达 $100\%$。可见窗口大小 $N$ 需满足 $N \ge \lceil (RTT + L/R)/(L/R) \rceil$ 才能完全利用带宽。

流水线带来了新问题：出错时如何处理？由此衍生出 GBN 与 SR 两种协议。

\subsection{回退 N 步 (GBN)}

GBN（Go-Back-N）中，发送方维护窗口 $[base,\; base+N-1]$：

\begin{itemize}
    \item \textbf{累计确认 (cumulative ACK)}：收到 ACK $n$ 表示序号 $n$ 及之前\emph{所有}分组都已正确收到，窗口起点滑到 $n+1$。
    \item \textbf{单一计时器}：只对最早的未确认分组计时。
    \item \textbf{超时即回退}：一旦超时或收到重复 ACK，\emph{重传从 $base$ 起的全部已发送分组}。
    \item 接收方只按序接收，\emph{丢弃乱序分组}，并对最后一个按序分组重复发 ACK。
\end{itemize}

GBN 接收窗口为 $1$，实现简单、所需缓存小，但一个分组出错会引起大量重传，浪费带宽。

\subsection{选择重传 (SR)}

SR（Selective Repeat）中，接收方为每个正确到达（即使乱序）的分组单独确认并缓存，发送方\emph{只重传真正丢失的分组}：

\begin{itemize}
    \item \textbf{逐分组确认}：每个分组有独立 ACK，不再累计。
    \item \textbf{多计时器}：每个未确认分组各自计时。
    \item \textbf{接收窗口 $>1$}：缓存乱序分组，待缺口补齐后一并按序交付。
\end{itemize}

SR 节省带宽、重传少，但需要更多缓存与更复杂的逻辑。

\begin{table}[H]
\centering
\caption{GBN 与 SR 的对比}
\begin{tabulary}{\textwidth}{CCC}
\toprule
维度 & GBN（回退 N 步） & SR（选择重传） \\
\midrule
确认方式 & 累计确认 & 逐分组确认 \\
计时器 & 仅最早未确认分组 & 每个未确认分组 \\
出错处理 & 重传 $base$ 起全部 & 只重传丢失分组 \\
接收缓存 & 不需要（窗口为 $1$） & 需要缓存乱序分组 \\
重传开销 & 大 & 小 \\
最大窗口 $W$ & $W \le 2^{k}-1$ & $W \le 2^{k-1}$ \\
\bottomrule
\end{tabulary}
\label{tab:gbn_sr}
\end{table}

\subsection{序号回绕与窗口上限}

序号用 $k$ 位表示，共有 $2^{k}$ 个不同序号，循环使用，称作\emph{序号回绕}。若序号空间太小、窗口太大，接收方可能把\emph{旧的重传分组}误认为\emph{新分组}。

\textbf{GBN 为何要求 $W \le 2^{k}-1$？} 假设 $k=2$（序号 $0,1,2,3$，空间 $2^2=4$），若窗口取 $W=4$。发送方发 $0,1,2,3$ 全部丢失，接收方什么都没收到。发送方超时重传 $0,1,2,3$。若其中「重传的 $0$」延迟到很久以后才到达，而此时接收方已经开始接收下一轮的新分组 $0,1,\dots$，就\emph{无法区分}这个 $0$ 是旧的重传还是新的分组。所以必须保证窗口小于序号空间：$W \le 2^{k}-1$。

\textbf{SR 为何要求 $W \le 2^{k-1}$？} SR 的接收窗口与发送窗口等大，接收方要缓存乱序分组。设 $W = 2^{k-1}$ 时，发送窗口和接收窗口占据的序号范围不重叠，尚能区分「新分组」与「重复分组」；一旦 $W > 2^{k-1}$，两个窗口的序号范围会重叠，接收方无法判断一个落在重叠区间的分组到底是新数据还是旧重复。故 SR 的上限是对 GBN 的一半。

\subsection{GBN 时序示例}

\begin{figure}[H]
\centering
\begin{tikzpicture}[font=\small, >=latex, thick]
    \draw (0,0) -- (0,-9.2) node[below, font=\small]{时间};
    \draw (5.6,0) -- (5.6,-9.2) node[below, font=\small]{时间};
    \node[font=\small\bfseries] at (0,0.45) {发送方};
    \node[font=\small\bfseries] at (5.6,0.45) {接收方};

    \draw[->] (0,-1.0) -- (5.6,-1.7) node[midway, above, font=\scriptsize]{pkt0};
    \draw[<-] (0,-2.1) -- (5.6,-1.8) node[midway, below, font=\scriptsize]{ACK0};

    \draw[->] (0,-2.6) -- (5.6,-3.3) node[midway, above, font=\scriptsize]{pkt1};
    \draw[<-] (0,-3.7) -- (5.6,-3.4) node[midway, below, font=\scriptsize]{ACK1};

    \draw[->] (0,-4.1) -- (2.9,-4.6) node[midway, above, font=\scriptsize]{pkt2};
    \node[font=\scriptsize, red, align=center] at (3.55,-4.85) {$\times$ 丢失};

    \draw[->] (0,-5.2) -- (5.6,-5.9) node[midway, above, font=\scriptsize]{pkt3};
    \node[font=\scriptsize, align=center] at (6.1,-6.2) {乱序丢弃};
    \draw[<-] (0,-6.5) -- (5.6,-6.2) node[midway, below, font=\scriptsize]{ACK1 (重复)};

    \draw[dashed, red] (0,-6.9) -- (5.6,-6.9) node[midway, above, font=\scriptsize]{超时};
    \draw[->] (0,-7.3) -- (5.6,-8.0) node[midway, above, font=\scriptsize]{pkt2 (重传)};
    \draw[->] (0,-8.2) -- (5.6,-8.9) node[midway, above, font=\scriptsize]{pkt3 (重传)};
    \draw[<-] (0,-9.1) -- (5.6,-8.9) node[midway, below, font=\scriptsize]{ACK3};
\end{tikzpicture}
\caption{GBN 时序：pkt2 丢失，发送方从 base 起回退重传 pkt2、pkt3}
\label{fig:gbn_timeline}
\end{figure}

\begin{algorithm}[H]
\caption{GBN 发送方规约}
\begin{algorithmic}[1]
\State \textbf{状态}: 窗口 $[base,\;base+N-1]$, $nextseqnum\gets base$
\Procedure{GbnSend}{$data$}
    \If{$nextseqnum < base + N$}
        \State 以序号 $nextseqnum$ 发送分组，并缓存其副本
        \State $nextseqnum \gets nextseqnum + 1$
    \Else
        \State 缓存数据，等待窗口腾出
    \EndIf
\EndProcedure
\State \textbf{收到 ACK $n$（累计）}: $base \gets n+1$
\State \textbf{超时}: 重传序号 $base$ 到 $nextseqnum-1$ 的全部分组，重启计时器
\end{algorithmic}
\end{algorithm}

\section{TCP：传输控制协议}

\subsection{TCP 的三大特性}

\begin{itemize}
    \item \textbf{面向连接}：通信前必须通过三次握手建立连接，通信后需四次挥手释放；连接是全双工的，双方可同时发送数据。
    \item \textbf{字节流}：TCP 把应用数据看成\emph{无结构的字节流}，不保留报文边界。发送方按需分段，接收方按字节序号重组；「粘包」正是字节流的体现。
    \item \textbf{全双工可靠}：两个方向各自独立确认、独立流量控制，配合序号、确认、超时重传与滑动窗口实现可靠传输。
\end{itemize}

TCP 一次发送的数据上限受 \emph{MSS}（最大报文段长度）约束，MSS 由链路 \emph{MTU}（最大传输单元，以太网 $1500$ 字节）减去 IP 与 TCP 首部得到（典型 \texttt{MSS}=$1460$ 字节）。

\subsection{序号与确认号语义}

TCP 首部中的两个关键字段是 \textbf{序号 (sequence number)} 和 \textbf{确认号 (acknowledgment number)}。理解它们的语义是读懂 TCP 的关键：

\begin{itemize}
    \item \textbf{序号}：本报文段\emph{数据部分第一个字节}在整个字节流中的编号（初始序号 ISN 随机选取，安全起见不固定为 $0$）。注意序号是\emph{字节序号}，不是报文序号。
    \item \textbf{确认号}：期望从对方收到的\emph{下一个字节}的序号，即「我已收到该字节之前的\emph{所有}字节」。这体现了 TCP 的\emph{累计确认}。
    \item \textbf{ACK 标志}：确认号字段有效当且仅当 ACK 标志置 $1$。
\end{itemize}

\textbf{算例}。客户机随机初始序号 $ISN_C = 1000$，已发送 $200$ 字节数据（字节 $1000\sim1199$）。服务器回复确认号 $1200$、序号 $ISN_S = 5000$，表示「我已收到 $1200$ 之前的全部字节，下一字节应为 $1200$」。此后客户机再发确认号 $5000$、序号 $1200$ 的报文段。

\subsection{连接管理}

\subsubsection{三次握手及「为什么不是两次」}

\begin{figure}[H]
\centering
\begin{tikzpicture}[font=\small, >=latex, thick]
    \node[font=\small\bfseries] at (0,0.5) {客户机};
    \node[font=\small\bfseries] at (6,0.5) {服务器};
    \draw (0,0) -- (0,-4.2);
    \draw (6,0) -- (6,-4.2);

    \draw[->] (0,-0.8) -- (6,-1.5) node[midway, above, align=center, font=\scriptsize]{SYN, seq=x};
    \node[font=\scriptsize, anchor=east] at (0,-0.8) {CLOSED};
    \node[font=\scriptsize, anchor=west] at (6,-1.5) {LISTEN};

    \draw[<-] (0,-2.2) -- (6,-1.6) node[midway, above, align=center, font=\scriptsize]{SYN+ACK, seq=y, ack=x+1};
    \node[font=\scriptsize, anchor=east] at (0,-2.2) {SYN\_SENT};
    \node[font=\scriptsize, anchor=west] at (6,-1.6) {SYN\_RCVD};

    \draw[->] (0,-3.0) -- (6,-3.7) node[midway, above, align=center, font=\scriptsize]{ACK, ack=y+1};
    \node[font=\scriptsize, anchor=east] at (0,-3.0) {ESTABLISHED};
    \node[font=\scriptsize, anchor=west] at (6,-3.7) {ESTABLISHED};
\end{tikzpicture}
\caption{TCP 三次握手}
\label{fig:three_way}
\end{figure}

\textbf{为什么两次握手不够？} 至少三条理由：

\begin{enumerate}
    \item \textbf{防止失效的旧连接请求}：若某次「客户机请求连接」的 SYN 因网络拥堵而\emph{滞留}，客户机超时后重发并完成了一次连接、通信、释放。此时那个滞留的旧 SYN 突然到达服务器。若只有两次握手，服务器收到后会立即认为连接建立并分配资源，而客户机根本不理它，导致服务器白白占用资源。三次握手时，服务器回 SYN+ACK 后\emph{等不到}客户机的第三次 ACK，就不会真正建立连接。
    \item \textbf{确认双方收发能力}：两次只能确认「客户机 $\to$ 服务器」方向可达；第三次 ACK 才能让服务器确认「服务器 $\to$ 客户机」方向也可达，即双方都确认了对方的收发能力。
    \item \textbf{同步双方初始序号}：TCP 可靠传输依赖字节序号，双方必须先交换并确认彼此的 ISN，三次握手恰好完成「请求、同意并交换序号、确认」三步。
\end{enumerate}

\subsubsection{四次挥手、TIME\_WAIT 与 2$\times$MSL}

\begin{figure}[H]
\centering
\begin{tikzpicture}[font=\small, >=latex, thick]
    \node[font=\small\bfseries] at (0,0.5) {主动关闭方};
    \node[font=\small\bfseries] at (6,0.5) {被动关闭方};
    \draw (0,0) -- (0,-4.6);
    \draw (6,0) -- (6,-4.6);

    \draw[->] (0,-0.8) -- (6,-1.4) node[midway, above, font=\scriptsize]{FIN, seq=u};
    \draw[<-] (0,-2.0) -- (6,-1.5) node[midway, above, font=\scriptsize]{ACK, ack=u+1};
    \node[font=\scriptsize, anchor=east] at (0,-2.0) {FIN\_WAIT\_2};

    \draw[<-] (0,-3.0) -- (6,-2.6) node[midway, above, font=\scriptsize]{FIN, seq=w};
    \draw[->] (0,-3.5) -- (6,-3.0) node[midway, above, font=\scriptsize]{ACK, ack=w+1};
    \node[font=\scriptsize, anchor=east, align=center] at (0,-4.0) {TIME\_WAIT\\$2\times$MSL};
    \node[font=\scriptsize, anchor=west] at (6,-3.0) {CLOSED};
\end{tikzpicture}
\caption{TCP 四次挥手（主动关闭方进入 TIME\_WAIT）}
\label{fig:four_way}
\end{figure}

TCP 连接是全双工的，每个方向需\emph{单独}关闭，因此是「FIN $\to$ ACK $\to$ FIN $\to$ ACK」四步。收到对方 FIN 后先回 ACK（此时可能还有未传完的数据要继续发送，故 ACK 与自己的 FIN 不能合并），待本方数据也发完，再发自己的 FIN。

主动关闭方在发出最后一个 ACK 后进入 \texttt{TIME\_WAIT} 状态，等待 $2\times$MSL（MSL 为报文最大生存时间，典型 $2$ 分钟）才彻底关闭。原因有二：

\begin{enumerate}
    \item \textbf{确保最后的 ACK 能到达}：若最后的 ACK 丢失，被动关闭方会超时重传 FIN；处于 \texttt{TIME\_WAIT} 的主动方还能再回一个 ACK。若立刻关闭，对方重传的 FIN 将无从应答，导致对方无法正常关闭。
    \item \textbf{让旧连接的报文在网络中消散}：等待 $2\times$MSL（一个往返的最大寿命）可确保本次连接的所有报文段都从网络中消失，不致干扰随后建立的、复用同一四元组的新连接。
\end{enumerate}

\subsection{RTT 估计与超时 (Jacobson/Karels)}

超时时间设得太短会引发不必要的重传，太长则对丢包反应迟钝。TCP 用采样与平滑估计 $RTT$。设某次实际测量得到 $\text{SampleRTT}$，则

\begin{align}
    \text{EstimatedRTT} &= (1-\alpha)\cdot\text{EstimatedRTT} + \alpha\cdot\text{SampleRTT}, \qquad \alpha = 0.125, \\
    \text{DevRTT} &= (1-\beta)\cdot\text{DevRTT} + \beta\cdot|\text{SampleRTT} - \text{EstimatedRTT}|, \qquad \beta = 0.25, \\
    \text{TimeoutInterval} &= \text{EstimatedRTT} + 4\cdot\text{DevRTT}.
    \label{eq:timeout}
\end{align}

\textbf{思想}：$\text{EstimatedRTT}$ 是加权移动平均，$\alpha=0.125$ 使新样本权重小、估计平滑；$\text{DevRTT}$ 度量 $RTT$ 的\emph{波动}（偏差）；超时值取「平滑值 $+$ 若干倍偏差」，在波动大时自动放宽，兼顾及时性与稳健性。

\textbf{算例}。初始 $\text{EstimatedRTT}=\text{SampleRTT}=100$ ms，$\text{DevRTT}=\text{SampleRTT}/2=50$ ms。随后依次收到两个样本：

\begin{tabulary}{\textwidth}{CCCC}
\toprule
新样本 & 更新后 EstimatedRTT & 更新后 DevRTT & 超时时间 \\
\midrule
$120$ ms & $0.875\times100+0.125\times120=102.5$ & $0.75\times50+0.25\times20=42.5$ & $102.5+4\times42.5=272.5$ \\
$90$ ms & $0.875\times102.5+0.125\times90=100.94$ & $0.75\times42.5+0.25\times10=34.38$ & $100.94+4\times34.38=238.4$ \\
\bottomrule
\end{tabulary}

可见随着 $RTT$ 波动减小，超时时间从 $300$ ms 逐步收紧到约 $238$ ms。

\subsection{快速重传}

仅靠超时重传，发送方要等到计时器到期，时延较大。\textbf{快速重传 (fast retransmit)} 的思想是：当发送方收到对\emph{同一个}字节的 \textbf{3 个重复 ACK} 时，基本可断定该分组已丢失（后续分组到达促使接收方反复发送同一累计确认），于是\emph{不必等超时}立即重传，显著缩短恢复时间。快速重传是快速恢复（见下节）的触发条件。

\section{流量控制与拥塞控制}

\subsection{流量控制：接收窗口 rwnd}

流量控制解决的是「\emph{发送方太快、接收方来不及收}」的问题。接收方在自己的 TCP 首部中用 \textbf{接收窗口 rwnd} 字段通告当前可用缓存的大小，发送方据此限制在途未确认字节数：

\begin{equation}
    \text{LastByteSent} - \text{LastByteAcked} \le \text{rwnd}.
    \label{eq:flow}
\end{equation}

当接收方缓存将满，rwnd 变小甚至为 $0$，发送方停止发送；接收方缓存腾空后再通告新的 rwnd。为避免 rwnd 变为 $0$ 后双方互相等待，TCP 还使用\emph{持续计时器}，让发送方周期性地发送探测报文询问窗口是否恢复。

\subsection{拥塞控制：拥塞窗口 cwnd}

拥塞控制解决的是「\emph{发送方太快、网络承受不了}」的问题。发送方额外维护一个 \textbf{拥塞窗口 cwnd}，反映网络当前可承受的发送量。实际允许发送的窗口为两者较小值：

\begin{equation}
    \text{允许发送窗口} = \min(\text{rwnd},\; \text{cwnd}).
    \label{eq:min_window}
\end{equation}

\subsection{慢启动、拥塞避免与快速恢复}

\begin{itemize}
    \item \textbf{慢启动 (slow start)}：cwnd 从 $1$ MSS 开始，每收到一个 ACK 加 $1$ MSS；一个 RTT 内翻一倍，呈\emph{指数增长}。当 cwnd 达到阈值 \texttt{ssthresh} 时转入拥塞避免。
    \item \textbf{拥塞避免 (congestion avoidance)}：每个 RTT 只把 cwnd 加 $1$ MSS，呈\emph{线性增长}，谨慎探测可用带宽。
    \item \textbf{快速恢复 (fast recovery)}：收到 $3$ 个重复 ACK 时，把 \texttt{ssthresh} 减半，cwnd 从新阈值继续（Reno）。
\end{itemize}

\begin{algorithm}[H]
\caption{TCP 拥塞控制（Reno）事件处理}
\begin{algorithmic}[1]
\State \textbf{初始}: $cwnd\gets1$ MSS, $ssthresh\gets$ 较大值
\For{每个 RTT}
    \If{$cwnd < ssthresh$} \Comment{慢启动}
        \State $cwnd \gets cwnd \times 2$
    \Else \Comment{拥塞避免}
        \State $cwnd \gets cwnd + 1$ MSS
    \EndIf
\EndFor
\State \textbf{收到 3 个重复 ACK}: $ssthresh\gets cwnd/2$;\; $cwnd\gets ssthresh$;\; 快速重传丢失分组 \Comment{进入快速恢复}
\State \textbf{超时}: $ssthresh\gets cwnd/2$;\; $cwnd\gets1$;\; 重新慢启动 \Comment{严重拥塞}
\end{algorithmic}
\end{algorithm}

\begin{figure}[H]
\centering
\begin{tikzpicture}[font=\small, >=latex, thick]
    \draw[->] (0,0) -- (7.5,0) node[right]{时间 / RTT};
    \draw[->] (0,0) -- (0,3.4) node[above]{cwnd};
    \draw[help lines] (0,2.2) -- (7.5,2.2) node[right, font=\scriptsize]{ssthresh};
    \draw[thick, blue] (0,0.25) -- (1,0.5) -- (2,1.0) -- (3,2.0) -- (3.3,2.2);
    \node[font=\scriptsize, blue, anchor=south east] at (2.0,1.0) {慢启动（指数）};
    \draw[thick, red] (3.3,2.2) -- (6.0,2.7);
    \node[font=\scriptsize, red, anchor=west] at (4.3,2.35) {拥塞避免（线性）};
    \draw[thick, red] (6.0,2.7) -- (6.0,1.35);
    \draw[thick, blue] (6.0,1.35) -- (7.4,1.9);
    \node[font=\scriptsize, anchor=west] at (6.05,1.15) {3 个重复 ACK: 减半};
\end{tikzpicture}
\caption{cwnd 演化：慢启动 $\to$ 拥塞避免 $\to$ 快速恢复}
\label{fig:cwnd}
\end{figure}

\subsection{Tahoe 与 Reno 的对比}

两个经典版本在检测到丢包后的反应不同：

\begin{table}[H]
\centering
\caption{TCP Tahoe 与 TCP Reno}
\begin{tabulary}{\textwidth}{CCC}
\toprule
触发事件 & Tahoe & Reno \\
\midrule
超时 & ssthresh 减半, cwnd$=1$, 慢启动 & ssthresh 减半, cwnd$=1$, 慢启动 \\
3 个重复 ACK & ssthresh 减半, cwnd$=1$, 慢启动 & 快速重传 $+$ 快速恢复, cwnd$=$ssthresh \\
是否区分两类丢包 & 否 & 是（超时更严重） \\
\bottomrule
\end{tabulary}
\label{tab:tahoe_reno}
\end{table}

Reno 比 Tahoe 更细：它认为「重复 ACK」说明后续分组仍在到达、网络尚有传输能力，因此只把窗口减半并快速恢复，而非直接退回 $1$。

\subsection{AIMD 与吞吐量近似}

拥塞控制的调整策略可概括为 \textbf{AIMD}（加性增、乘性减）：

\begin{itemize}
    \item \textbf{加性增 (Additive Increase)}：无拥塞时每个 RTT 窗口 $+1$ MSS，线性上升，缓慢试探。
    \item \textbf{乘性减 (Multiplicative Decrease)}：检测到丢包时窗口\emph{减半}，快速退出拥塞。
\end{itemize}

AIMD 使窗口呈\emph{锯齿波 (sawtooth)}：从 $W/2$ 线性爬到 $W$，遇丢包跌回 $W/2$，如此往复。窗口在一个周期内的平均值为

\[
\bar{W} = \frac{W/2 + W}{2} = \frac{3}{4}W = 0.75\,W.
\]

因此平均吞吐量近似为

\begin{equation}
    \text{吞吐量} \approx \frac{\bar{W}}{RTT} = \frac{0.75\,W}{RTT},
    \label{eq:throughput}
\end{equation}

其中 $W$ 为丢包发生时的窗口大小。上式说明：吞吐量与窗口成正比、与 $RTT$ 成反比——长肥管道（高 $RTT$、高带宽）中的单个 TCP 连接很难跑满带宽，这正推动了多连接并行与后续 TCP 变体的发展。

\subsection{TCP 公平性}

若 $k$ 条 TCP 连接共享一条带宽为 $R$ 的瓶颈链路，理想情况下每条连接最终获得

\begin{equation}
    \text{每条连接吞吐} \approx \frac{R}{k},
    \label{eq:fairness}
\end{equation}

即公平地均分带宽。其直观解释是 AIMD 的收敛性：窗口较大的连接增长同样快、但在拥塞时减半损失也更大，两条连接的「工作点」会逐渐逼近等分线。当然，若连接 $RTT$ 差异大或应用开多条并行连接，公平性会被破坏，这也是 TCP 公平性讨论的边界。

\begin{figure}[H]
\centering
\begin{tikzpicture}[font=\small, >=latex, thick,
    ax/.style={->, thick}]
    \draw[ax] (0,0) -- (5.2,0) node[right]{连接 1 吞吐 $x_1$};
    \draw[ax] (0,0) -- (0,5.2) node[above]{连接 2 吞吐 $x_2$};
    \draw[dashed] (0,0) -- (4.6,4.6) node[above right, font=\scriptsize]{公平线 $x_1=x_2$};
    \draw[help lines] (4.6,0) -- (0,4.6);
    \coordinate (p) at (3.6,1.4);
    \fill (p) circle (2.2pt);
    \draw[->, blue] (p) -- ++(0.55,0.55) node[midway, above left, font=\scriptsize]{$+1$ (AI)};
    \draw[->, red] (p) -- ++(-1.0,0) node[midway, below, font=\scriptsize]{$\div 2$ (MD)};
    \node[font=\scriptsize, anchor=west] at (3.7,1.5) {当前工作点};
\end{tikzpicture}
\caption{AIMD 的向量分解：加性增沿 $45^\circ$ 逼近公平线，乘性减向原点收缩}
\label{fig:fairness}
\end{figure}

\section{本章小结}

\begin{itemize}
    \item 运输层在\emph{进程到进程}层面提供逻辑通信；发送端\emph{复用}、接收端\emph{分用}，依据是端口号。
    \item \textbf{UDP 用二元组} (目的 IP, 目的端口) 分用，因其无连接、无状态；\textbf{TCP 用四元组}分用，因每条连接需要独立状态。
    \item UDP 首部 $8$ 字节，检验和用 $16$ 位反码和覆盖伪首部、首部与数据；适合时延敏感、可容忍丢包或需自行控制的应用。
    \item 可靠传输由确认、序号、超时、窗口组合实现。停等利用率 $U=\frac{L/R}{RTT+L/R}$，滑动窗口 $U=\frac{N L/R}{RTT+L/R}$。GBN 用累计确认、窗口 $\le 2^{k}-1$；SR 逐分组确认、窗口 $\le 2^{k-1}$。
    \item TCP 面向连接、字节流、全双工；序号与确认号均为\emph{字节}语义且累计确认；三次握手防旧连接并同步初始序号，四次挥手因双向各需关闭，主动方 \texttt{TIME\_WAIT} 等 $2\times$MSL。
    \item 超时用 Jacobson/Karels 公式 $\text{TimeoutInterval}=\text{EstimatedRTT}+4\cdot\text{DevRTT}$；$3$ 个重复 ACK 触发快速重传。
    \item 流量控制看 \emph{rwnd}（接收方能力），拥塞控制看 \emph{cwnd}（网络能力），实际窗口为 $\min(\text{rwnd},\text{cwnd})$。慢启动指数增、拥塞避免线性增、快速恢复；Tahoe 一律退回慢启动，Reno 用快速恢复；AIMD 使吞吐近似 $\frac{0.75W}{RTT}$ 并趋向公平。
\end{itemize}

\subsection*{常见误区}

\begin{enumerate}
    \item \textbf{混淆流量控制与拥塞控制}：流量控制是\emph{接收方}通告 rwnd 以免被淹没；拥塞控制是\emph{网络}拥塞时减小 cwnd。二者共同决定发送窗口 $\min(\text{rwnd},\text{cwnd})$，但目的与触发源不同。
    \item \textbf{把 TCP 序号当成报文序号}：TCP 序号是\emph{字节}编号，确认号是「期望的下一个字节」，两者都以字节为单位。
    \item \textbf{认为两次握手也够用}：两次无法防止失效的旧 SYN，也无法让服务器确认自己的发送方向可达。
    \item \textbf{认为 TIME\_WAIT 是多余的等待}：它用于重发可能丢失的最后 ACK，并让旧连接报文在网络中消散。
    \item \textbf{把 GBN 与 SR 的窗口上限记反}：GBN 为 $2^{k}-1$，SR 为 $2^{k-1}$。
    \item \textbf{误以为 UDP 没有校验或一定不可靠到不可用}：UDP 有校验和（IPv4 中可选）；其可靠性可由应用层补足。
\end{enumerate}
